\section{Extension to independent cell-local controls}
\label{app:local-controls}

We first connect a microscopic control to the enlarged probe space in
Eq.~\eqref{eq:local-control-probe}. Add
$\delta H_{\rm loc}=\beta\sum_i W_i$, where each bounded Hermitian $W_i$
is an SU(2) scalar confined to cell $i$. The reference ground multiplet is
an isolated single spin-$s$ irrep. Schur's lemma gives
$P_iW_iP_i=w_iP_i$, a constant rather than an effective spin interaction.
The first-order derivative of the local vector probe is again a vector.
Because rank one occurs once in the operator space of a single spin-$s$
irrep, its projection must be a multiple of $T_i^z$. At $q=\pi$ the
probe is Hermitian, so this multiple is real. Changes in the projected
end-spin coefficients enter $\beta V$ only at order $\beta^2$;
likewise, the change of $b$ affects the physical first-order source only
at order $\beta^2$. The leading Hamiltonian is consequently still the
same uniform $h$, and all such controls are contained in the real
one-spin space of Eq.~\eqref{eq:local-control-probe}. Cell reflection
symmetry need not be preserved by $W_i$. This is a tangent-space result
about the reference root, not a finite displacement along a family of
different microscopic models.

Let $U$ translate by one effective site. For eigenstates of $h$, define
the positive semidefinite sesquilinear form
\begin{equation}
 (A,B)_{\tau,\nu}=\sum_{\epsilon_f-\epsilon_i>\nu}
 p_i A_{fi}^*B_{fi},\qquad \nu\ge0,
 \label{eq:window-inner-product}
\end{equation}
where $p_i$ are Gibbs weights or $1/d^N$. All degeneracies are included
as complete blocks. Since $[U,h]=0$, conjugation by $U$ is an isometry
of this form. Distinct operator momenta are orthogonal even when the
form has a null space. The $\pi$-momentum projection is
\begin{equation}
 \Pi_\pi A=\frac1N\sum_{m=0}^{N-1}(-1)^mU^m A U^{-m}.
\end{equation}
For even $N$, it obeys
$\Pi_\pi R_{\{u_i\}}=R_{\bar u}$. Pythagoras in
Eq.~\eqref{eq:window-inner-product} gives
Eq.~\eqref{eq:local-control-projection}, and hence
\begin{equation}
 \min_{\{u_i\}\in\mathbb R^N}F(R_{\{u_i\}})
 =\min_{u\in\mathbb R}F(R_u).
 \label{eq:local-control-optimum}
\end{equation}
The equality uses the availability of uniform controls in the enlarged
space. Any more restricted microscopic control set still satisfies the
same lower bound. The nonzero derivative $A'$ makes the uniform direction
accessible in the pentamer example. This argument applies at every
positive scaled temperature and to the uniform-trace window bound;
it does not require a nondegenerate spectrum.

The second moment provides an additional explicit check. On bond
$(i,i+1)$ the dipole commutator coefficient becomes
$(u_i+u_{i+1})/2$. Repeating the support-orthogonality calculation of
Appendix~\ref{app:general-moment} gives
\begin{align}
 M_2(\{u_i\})={}&M_{2,\min}+\frac{(d^2-1)^2}{18\ell N}\notag\\
 &\times\sum_i\left[\frac{u_i+u_{i+1}}2+
              \frac{b(d^2-4)}{12}\right]^2.
 \label{eq:nonuniform-moment}
\end{align}
An alternating component of $u_i$ adds a conserved total $T^z$ to the
probe and leaves this moment unchanged. For the same reason, the earlier
$M_2/M_0$ inequality cannot be applied unchanged to an arbitrary
nonuniform control. The finite-frequency lower bound follows instead
from the positive operator-momentum decomposition above.

For all ten pairwise exchange directions inside the isolated pentamer,
the derivative $-P(W_iGB+BGW_i)P$, with
$G=Q(H_{\rm cell}-E_0)^{-1}Q$, has a non-dipole Frobenius norm below
$2.4\times10^{-15}$. Central differences at step $5\times10^{-5}$
agree within $2.8\times10^{-10}$. Full four-site spectra for $s=1,3/2$
at $\tau=0.2,1.6$ and uniform trace verify
Eq.~\eqref{eq:local-control-optimum} for $\nu=0$ and
$\sqrt{3s(s+1)/4}$, with maximum difference $1.7\times10^{-16}$ at
$b=1$. Full-chain matrices at $N=4,6$ verify
Eq.~\eqref{eq:nonuniform-moment} within $1.5\times10^{-14}$.
